% \iffalse meta-comment
% File: sbornik-font.dtx (part of the sbornik bundle)
% \fi
%
% \iffalse
%<*driver>
\documentclass[full]{l3doc}
\begin{document}\DocInput{\jobname.dtx}
\end{document}
%</driver>
% \fi
%
% \title{Module \pkg{sbornik-font}}
% \author{}\date{}\maketitle
%
% \begin{documentation}
% Font-size module. В классе регистрирует ключи \texttt{fontsize=NNpt}
% и «голые» \texttt{NNpt}.
% \end{documentation}
%
% \begin{implementation}
%    \begin{macrocode}
%<*module>
\tl_new:N \g_sbornik_fontsize_tl
\tl_gset:Nn \g_sbornik_fontsize_tl { 10pt }

\keys_define:nn { sbornik/font }
{
  fontsize .tl_gset:N = \g_sbornik_fontsize_tl,
  fontsize .initial:n = 10pt,
}

% Регистрирует в семействе класса ключи fontsize и «голые» NNpt.
\cs_new_protected:Npn \sbornik_font_declare_keys:n #1
{
  \keys_define:nn {#1}
  {
    fontsize .code:n =
    { \keys_set:nn { sbornik/font } { fontsize = ##1 } },
    8pt  .meta:n = { fontsize = 8pt },
    9pt  .meta:n = { fontsize = 9pt },
    10pt .meta:n = { fontsize = 10pt },
    11pt .meta:n = { fontsize = 11pt },
    12pt .meta:n = { fontsize = 12pt },
    14pt .meta:n = { fontsize = 14pt },
    17pt .meta:n = { fontsize = 17pt },
    20pt .meta:n = { fontsize = 20pt },
  }
}
%</module>
%    \end{macrocode}
% \end{implementation}
% \Finale
\endinput
